
\SGASection{4}{Characterization of a sub-torus \(T\) by the subgroups
\({}_nT\)}
\SGARefTarget{sga3:xv:section:4}

\SGASubsection{4.1}{Statement of the main theorem}
\SGARefTarget{sga3:xv:subsection:4:1}

\SGAStatement{\SGAInlineRefTarget{sga3:xv:theorem:4-1}Theorem 4.1}{Let \(S\) be a connected locally
noetherian prescheme, let \(G\) be an \(S\)-group prescheme of finite
type over \(S\), let \(q>1\) be an integer invertible on \(S\), and
let \(r\) be a positive integer. For every integer \(n\) that is a
power of \(q\), let \(M(n)\) be a subgroup scheme of \(G\), of
multiplicative type and of type
\((\mathbb Z/n\mathbb Z)^r\). Suppose that:
\begin{enumerate}
\renewcommand{\labelenumi}{\textup{\alph{enumi})}}
\item the family \(M(n)\) is coherent: if \(m\mid n\), then
\[
  {}_mM(n)=M(m);
\]
\item\SGAInlineRefTarget{sga3:xv:theorem:4-1:item:b} there exist a point \(s\in S\) and a sub-torus \(T_s\) of
\(G_s\) such that
\[
  M(n)_s={}_nT_s\qquad\text{for every }n;
\]
\item for every point \(t\in S\), there exists a closed affine
subscheme \(F_t\) of \(G_t\) containing \(M(n)_t\) for every \(n\).
\end{enumerate}
Then there exists a unique sub-torus \(T\) of \(G\) such that
\({}_nT=M(n)\) for every \(n\) that is a power of \(q\).}
\SGARefTarget{sga3:xv:theorem:4:1}

There is an analogous theorem for lifting morphisms.

\SGAStatement{\SGAInlineRefTarget{sga3:xv:theorem:4-1-bis}Theorem 4.1 bis}{Let \(S,G,q\) be as above, let \(T\)
be an \(S\)-torus, and, for every \(n\) that is a power of \(q\), let
\(u(n):{}_nT\to G\) be an \(S\)-homomorphism. Suppose that:
\begin{enumerate}
\renewcommand{\labelenumi}{\textup{\alph{enumi})}}
\item the family \(u(n)\) is coherent: if \(m\mid n\), then
\[
  u(m)=u(n)|_{{}_mT};
\]
\item there exist a point \(s\in S\) and a homomorphism
\(u_s:T_s\to G_s\) such that
\[
  u_s|_{{}_nT_s}=u(n)_s
  \qquad\text{for every }n\text{ that is a power of }q;
\]
\item for every point \(t\in S\), there exists a closed affine
subscheme \(F_t\) of \(G_t\) containing
\(u(n)_t({}_nT_t)\) for every \(n\).
\end{enumerate}
Then there exists a unique homomorphism \(u:T\to G\) whose restriction
to \({}_nT\) equals \(u(n)\) for every \(n\) that is a power of \(q\).}
\SGARefTarget{sga3:xv:theorem:4:1:bis}

\SGAStatement{Remark 4.2}{Using the lower semicontinuity of the
abelian rank of a flat group prescheme of finite type over the
spectrum of a discrete valuation ring (cf. Exposé \SGARefLink{sga3:x:remark:8-7}{X~8.7}), condition
(c) in \SGARefLink{sga3:xv:theorem:4-1}{Theorems 4.1} and \SGARefLink{sga3:xv:theorem:4-1-bis}{4.1 bis} may be weakened by requiring the
closed affine subscheme \(F_t\) only for every maximal point \(t\) of
\(S\).}
\SGARefTarget{sga3:xv:remark:4:2}

Let us deduce \SGARefLink{sga3:xv:theorem:4-1-bis}{Theorem 4.1 bis} from \SGARefLink{sga3:xv:theorem:4-1}{Theorem 4.1}. Put
\(G'=G\times_S T\). For every \(n\) that is a power of \(q\), consider
the homomorphism
\[
  v(n):{}_nT\longrightarrow G'
\]
whose projections to \(G\) and \(T\) are respectively \(u(n)\) and
the canonical immersion \({}_nT\to T\). Thus \(v(n)\) is an immersion;
let \(M(n)\) be its image subgroup. The family \(M(n)\) is coherent
in the sense of \SGARefLink{sga3:xv:theorem:4-1}{Theorem 4.1}. The group \(M(n)_s\) equals
\({}_nT'_s\), where \(T'_s\) is the sub-torus of \(G'_s\) that is the
graph of \(u_s\), and, for every \(t\in S\), the closed affine
subscheme \(F_t\times_tT_t\) of \(G'_t\) contains all the groups
\(M(n)_t\). By \SGARefLink{sga3:xv:theorem:4-1}{Theorem 4.1} there is therefore a sub-torus \(T'\) of
\(G'\) such that \({}_nT'=M(n)\) for every \(n\).

Let \(f:T'\to T\) be the restriction of the projection \(G'\to T\),
and let \(f(n)\) be its restriction to \(M(n)\). Then
\[
  f(n)\circ v(n)=\operatorname{id}_{{}_nT}.
\]
The fibre \(T'_s\) is the graph of \(u_s\) (Exposé \SGARefLink{sga3:ix:corollary:4:8:item:b}{IX~4.8(b)}), so
\(f_s\) is an isomorphism. The groups \(\ker f\) and
\(\operatorname{coker}f\) are of multiplicative type (Exposé \SGARefLink{sga3:ix:proposition:2:7}{IX~2.7})
and of constant type, since \(S\) is connected. They are therefore
trivial, and \(f\) is an isomorphism. If \(v=f^{-1}\), then
\[
  v|_{{}_nT}=v(n).
\]
Composing \(v\) with the projection \(G'\to G\) gives the required
homomorphism \(u:T\to G\). This proves existence; uniqueness follows
from Exposé \SGARefLink{sga3:ix:corollary:4:8:item:a}{IX~4.8(a)}.

\SGASubsection{4.2}{Application}
\SGARefTarget{sga3:xv:subsection:4:2}

We shall generalize Theorem \SGARefLink{sga3:ix:theorem:7:1}{IX~7.1}. Let \(A\) be a complete
noetherian local ring, let \(I\) be its maximal ideal, and put
\[
  S=\operatorname{Spec}A,\qquad
  S_m=\operatorname{Spec}(A/I^m).
\]
For every prescheme \(X\), put \(X_m=X\times_SS_m\).

Let \(G\) be an \(S\)-group prescheme of finite type, let \(T\) be an
\(S\)-torus, let \(q\) be an integer invertible on \(S\), and let
\[
  u_m:T_m\longrightarrow G_m\qquad(m\in\mathbb N)
\]
be a coherent family of homomorphisms. For \(n\) ranging over the
powers of \(q\), denote by \(u_m(n)\) the restriction of \(u_m\) to
\({}_nT_m\), and by
\[
  u(n):{}_nT\longrightarrow G
\]
the unique homomorphism extending the morphisms \(u_m(n)\) for all
\(m\) (\SGARefLink{sga3:xv:proposition:1-6:item:a}{Proposition 1.6(a)}). We call the family \((u_m)_{m\in\mathbb N}\)
\emph{admissible} if, for every point \(t\in S\), there exists a
closed affine subprescheme \(F_t\) of \(G_t\) containing
\(u(n)_t({}_nT_t)\) for every \(n\) that is a power of \(q\). As will
be seen, this property does not depend on \(q\).

\SGAStatement{Proposition 4.3}{With the preceding notation, the
canonical map
\[
 \operatorname{Hom}_{S\text{-}\mathrm{gr}}(T,G)
 \longrightarrow
 \varprojlim_m\operatorname{Hom}_{S_m\text{-}\mathrm{gr}}(T_m,G_m)
\]
identifies its source with the subset of the target consisting of
admissible coherent families.}
\SGARefTarget{sga3:xv:proposition:4:3}

Indeed, apply \SGARefLink{sga3:xv:theorem:4-1-bis}{Theorem 4.1 bis} with \(s\) the closed point of \(S\).

\SGAStatement{Corollary 4.4}{With the preceding notation, suppose in
addition that \(G\) has affine fibres. Then the canonical map
\[
 \operatorname{Hom}_{S\text{-}\mathrm{gr}}(T,G)
 \longrightarrow
 \varprojlim_m\operatorname{Hom}_{S_m\text{-}\mathrm{gr}}(T_m,G_m)
\]
is an isomorphism.}
\SGARefTarget{sga3:xv:corollary:4:4}

Indeed, when \(G\) has affine fibres every coherent family is
admissible.

\SGAStatement{Remark 4.5}{When \(G\) is separated, the ring \(A\) in
Proposition 4.3 may be replaced by a noetherian \(I\)-adic ring: one
uses EGA III~5.4.1 in place of Proposition 1.6(a).}
\SGARefTarget{sga3:xv:remark:4:5}

\SGASubsection{4.3}{Proof of \SGARefLink{sga3:xv:theorem:4-1}{Theorem 4.1}}
\SGARefTarget{sga3:xv:subsection:4:3}

\SGAStatement{Lemma 4.6}{Let \(k\) be a field, let \(G\) be a
\(k\)-algebraic group, let \(q\) be an integer prime to the
characteristic of \(k\), let \(r>0\), and let \(M(n)\), for \(n\)
ranging over the powers of \(q\), be a coherent family of subgroups
of \(G\), of multiplicative type and of type
\((\mathbb Z/n\mathbb Z)^r\). There exists a smallest closed
subscheme \(H\) of \(G\) containing all the \(M(n)\). Moreover, \(H\)
is a smooth, connected, commutative algebraic subgroup of \(G\), and
its formation commutes with extension of the base field.}
\SGARefTarget{sga3:xv:lemma:4:6}

Let \(M\subseteq\operatorname{ens}(G)\) be the union of the underlying
sets of all the \(M(n)\), and let \(H\) be the reduced closed
subscheme of \(G\) whose underlying space is the closure of \(M\).
Each \(M(n)\) is étale over \(k\), hence reduced, and is therefore
contained in \(H\). Thus \(H\) is the smallest closed subscheme
containing all the \(M(n)\).

Let \(\bar k\) be an algebraically closed extension of \(k\). By
construction the \(M(n)\) are schematically dense in \(H\) (Exposé
\SGARefLink{sga3:ix:definition:4:1}{IX~4.1}), hence the \(M(n)_{\bar k}\) are schematically dense in
\(H_{\bar k}\) (Exposé \SGARefLink{sga3:ix:lemma:4:5}{IX~4.5}). Since every \(M(n)_{\bar k}\) is
reduced, \(H_{\bar k}\) is the reduced closed subscheme of
\(G_{\bar k}\) whose underlying space is the closure of \(M_{\bar k}\).
Thus \(H\) is geometrically reduced (EGA IV~4.6.1).

Coherence of the family implies that \(M\) is stable under the group
law. The subset \(M\times_kM\) is dense in
\(\operatorname{ens}(H\times_kH)\), and \(H\times_kH\) is reduced
(EGA IV~4.6.5); it follows that \(H\) is an algebraic subgroup of
\(G\). It is smooth over \(k\), since it is geometrically reduced
(Exposé VI\textsubscript{A}~\SGARefLink{prop:sga3-VIA-1.3.1}{1.3.1}), and it is commutative because all
the \(M(n)\) are commutative.

It remains to prove connectedness. Let \(H^0\) be the identity
component of \(H\), let \(m\) be the number of geometric points of
\(H/H^0\), and let \(q^s\) be the \(q\)-primary part of \(m\). For
every \(n\) that is a power of \(q\), the subgroup
\({}_{q^s}M(n)\) is contained in \(H^0\). Coherence and the prescribed
type of \(M(n)\) give
\[
  {}_{q^s}M(nq^s)=M(n).
\]
Hence \(H^0\) contains every \(M(n)\), and therefore \(H^0=H\).

We can now sharpen condition (c) of \SGARefLink{sga3:xv:theorem:4-1}{Theorem 4.1}. By \SGARefLink{sga3:xv:lemma:4:6}{Lemma 4.6} there
is, for each \(t\in S\), a smallest closed subscheme \(H_t\) of
\(G_t\) containing all the \(M(n)_t\). Its formation commutes with
extension of the residue field. Thus \(H_t\otimes_t\bar t\) is
contained in the affine closed subset \(F_t\), and is affine; hence
\(H_t\) itself is affine. \SGARefLink{sga3:xv:lemma:4:6}{Lemma 4.6} also says that \(H_t\) is a
smooth, connected, commutative algebraic group. The structure theorem
for such affine groups over an algebraically closed field
(\emph{Bible}, 4, \SGARefLink{sga3:xv:section:4}{Theorem 4}) expresses \(H_t\) as a direct product of
a torus \(T_t\) and a unipotent group \(U_t\). Since every \(M(n)_t\)
is contained in \(T_t\), minimality gives \(H_t=T_t\). Consequently
\(H_t\) is a torus.

We begin the proof of \SGARefLink{sga3:xv:theorem:4-1}{Theorem 4.1}.

\medskip\noindent
\textbf{a) Uniqueness.}
This is Exposé \SGARefLink{sga3:ix:corollary:4:8:item:b}{IX~4.8(b)}.

\medskip\noindent
\textbf{b) The case in which \(S\) is the spectrum of a complete
discrete valuation ring \(A\) with algebraically closed residue field
\(k\).}
Let \(K\) be the fraction field of \(A\), let \(s\) and \(t\) be the
closed and generic points of \(S\), let \(J\) be the maximal ideal of
\(A\), and put
\[
  S_m=\operatorname{Spec}(A/J^m),\qquad X_m=X\times_SS_m.
\]

\smallskip\noindent
\textbf{First case: the point in \SGARefLink{sga3:xv:theorem:4-1:item:b}{Theorem 4.1(b)} is the generic point
\(t\).}
Let \(T'\) be the schematic closure of \(T_t\) in \(G\). Its closed
fibre \(T'_s\) is an algebraic subgroup of \(G_s\), of dimension
\(r\), containing every \(M(n)_s\), hence containing \(H_s\). The
torus \(H_s\) contains every \(M(n)_s\), so it has rank at least
\(r\). Consequently \(H_s\) has the same underlying space as the
identity component of \(T'_s\).

The identity component \((T')^0\) of \(T'\) is then a subgroup
prescheme of \(G\), flat and separated (Exposé
VI\textsubscript{B}~\SGARefLink{sga3:xv:theorem:5-2}{5.2}), whose generic fibre \(T_t\) is a torus and
whose reduced closed fibre \(H_s\) is a torus. Hence \((T')^0\) is a
torus (Exposé \SGARefLink{sga3:x:theorem:8-8}{X~8.8}); denote it by \(T\). The groups \({}_nT\) and
\(M(n)\) are smooth over \(S\), hence reduced, and have the same
underlying space. They therefore coincide, and \(T\) is the required
solution.

\smallskip\noindent
\textbf{Second case: the point in \SGARefLink{sga3:xv:theorem:4-1:item:b}{Theorem 4.1(b)} is the closed point
\(s\).}
After replacing \(G\) by the schematic closure in \(G\) of the
smallest algebraic subgroup \(H_t\) containing the family \(M(n)_t\)
(\SGARefLink{sga3:xv:lemma:4:6}{Lemma 4.6}), we may suppose that \(G_t\) is affine.

For every \(m\geq0\), \SGARefLink{sga3:xv:subsection:2:2}{Proposition 2.2} gives a unique sub-torus \(T_m\)
of \(G_m\) lifting \(T_s\) and satisfying
\[
  {}_nT_m=M(n)_m
  \qquad\text{for every }n\text{ that is a power of }q.
\]
Put \(T'=\mathbb G_{m,S}^{\,r}\). Since \(k\) is algebraically
closed, \(T_s\) is split, and there is a \(k\)-isomorphism
\[
  u_s:T'_s\xrightarrow{\sim}T_s.
\]
It lifts uniquely to an \(S_m\)-isomorphism
\(u_m:T'_m\to T_m\) (Exposé \SGARefLink{sga3:ix:corollary:3:3}{IX~3.3}). The coherent family
\((u_m)_{m\in\mathbb N}\) defines a morphism of formal completions
\(\widehat u:\widehat{T'}\to\widehat G\), fitting into
\[
\begin{tikzcd}[column sep=large,row sep=large]
\widehat{T'} \arrow[r,"\widehat u"] \arrow[d,"i"'] &
  \widehat G \arrow[d,"j"'] \\
T' & G ,
\end{tikzcd}
\]
We shall prove that \(\widehat u\) is algebraizable, reducing first
to the case in which \(G\) is affine.

\SGAStatement{Lemma 4.7}{Let \(S=\operatorname{Spec}A\), where \(A\)
is a discrete valuation ring, and let \(X,Y\) be \(S\)-preschemes
that are quasi-compact, quasi-separated, and flat over \(S\). Then the
canonical map
\[
 \Gamma(X,\mathcal O_X)\otimes_A\Gamma(Y,\mathcal O_Y)
 \longrightarrow
 \Gamma(X\times_SY,\mathcal O_{X\times_SY})
\]
is an isomorphism.}
\SGARefTarget{sga3:xv:lemma:4:7}

Let \(f:X\to S\) and \(g:Y\to S\) be the structural morphisms.
Because \(X\) and \(Y\) are quasi-compact and quasi-separated,
EGA I~9.2.2 and EGA IV~1.7.4 show that
\(f_*\mathcal O_X\) and \(g_*\mathcal O_Y\) are quasi-coherent
\(\mathcal O_S\)-algebras. Let \(\widetilde X\) and \(\widetilde Y\)
be the corresponding affine \(S\)-schemes. Since \(Y\) is flat over
\(S\), EGA III~1.4.15 and EGA IV~1.7.21 give an isomorphism
\[
 \Gamma(\widetilde X\times_SY,
        \mathcal O_{\widetilde X\times_SY})
 \xrightarrow{\sim}
 \Gamma(X\times_SY,\mathcal O_{X\times_SY}).
\]
The scheme \(\widetilde X\) is flat over \(S\), because over \(A\)
flatness is equivalent to torsion-freeness. Applying EGA III~1.4.15
again, with \(X\) and \(Y\) interchanged, yields
\[
\begin{aligned}
 \Gamma(\widetilde X,\mathcal O_{\widetilde X})
 \otimes_A\Gamma(\widetilde Y,\mathcal O_{\widetilde Y})
 &\xrightarrow{\sim}
 \Gamma(\widetilde X\times_S\widetilde Y,
        \mathcal O_{\widetilde X\times_S\widetilde Y})\\
 &\xrightarrow{\sim}
 \Gamma(\widetilde X\times_SY,
        \mathcal O_{\widetilde X\times_SY}),
\end{aligned}
\]
which proves the lemma.

The \(S\)-group \(G\) is flat and of finite type over \(S\), hence
quasi-compact and quasi-separated. Applying the lemma to
\(G\times_SG\) gives
\[
 \Gamma(G,\mathcal O_G)\otimes_A\Gamma(G,\mathcal O_G)
 \xrightarrow{\sim}
 \Gamma(G\times_SG,\mathcal O_{G\times_SG}).
\]
Thus the multiplication \(m_G:G\times_SG\to G\) determines a
homomorphism
\[
 \Gamma(G,\mathcal O_G)\longrightarrow
 \Gamma(G,\mathcal O_G)\otimes_A\Gamma(G,\mathcal O_G),
\]
and hence an \(S\)-morphism
\[
  m_{\widetilde G}:\widetilde G\times_S\widetilde G
  \longrightarrow\widetilde G,
\]
where \(\widetilde G\) is the affine \(S\)-scheme corresponding to
\(\Gamma(G,\mathcal O_G)\). The usual identities show that
\(m_{\widetilde G}\) makes \(\widetilde G\) an \(S\)-group scheme and
that the canonical morphism \(v:G\to\widetilde G\) is a homomorphism.

\SGAStatement{Remark 4.8}{The group \(\widetilde G\) plays the role of
a largest affine quotient of \(G\). In fact it is an fpqc quotient of
\(G\), and is therefore of finite type over \(S\) (XVII, Appendix
III, 2).}
\SGARefTarget{sga3:xv:remark:4:8}

In our situation the generic fibre \(G_t\) is affine, so
\(v_t:G_t\to\widetilde G_t\) is an isomorphism (EGA I~9.3.3). Since
\(\widetilde G\) is affine, the coherent family
\[
  w_m=v_m\circ u_m:T'_m\longrightarrow\widetilde G_m
  \qquad(m\in\mathbb N)
\]
comes from a unique \(S\)-homomorphism (Exposé \SGARefLink{sga3:ix:theorem:7:1}{IX~7.1})
\[
  w:T'\longrightarrow\widetilde G.
\]
Let \(T_t\) be the sub-torus
\[
  T_t=v_t^{-1}\bigl(w_t(T'_t)\bigr)
\]
of \(G_t\). Its rank is at most \(r\). We show that it contains every
\(M(n)_t\). Let
\[
 u(n)_m:({}_nT')_m\xrightarrow{\sim}M(n)_m
\]
be the restriction of \(u_m\). This coherent family comes from a
unique \(S\)-isomorphism
\[
  u(n):{}_nT'\xrightarrow{\sim}M(n),
\]
because \(M(n)\) is finite over \(S\). For every \(m\geq0\),
\[
 w_m|_{({}_nT')_m}
  =(v_m\circ u_m)|_{({}_nT')_m}
  =(v\circ u(n))_m.
\]
It follows from \SGARefLink{sga3:xv:proposition:1-6:item:a}{Proposition 1.6(a)} that
\[
  w|_{{}_nT'}=v\circ u(n).
\]
On the generic fibre,
\[
 v_t^{-1}\!\left(w_t(({}_nT')_t)\right)
 =u(n)_t(({}_nT')_t)=M(n)_t.
\]
Thus \(T_t\) contains every \(M(n)_t\), and consequently has rank
\(r\).

We conclude as in the first case. Let \(T''\) be the schematic closure
of \(T_t\) in \(G\). Since \(T_t\) contains every \(M(n)_t\), the
closed fibre \(T''_s\) contains every \(M(n)_s\), and hence contains
\(T_s\) by the density theorem. On the other hand \(T''\) is flat
over \(S\) and \(T_t\) has dimension \(r\), so \(T''_s\) has
dimension \(r\) (Exposé VI\textsubscript{B}~4). Therefore \(T_s\)
has the same underlying space as the identity component of \(T''_s\),
and the identity component of \(T''\) is the required sub-torus of
\(G\).

\medskip\noindent
\textbf{c) End of the proof.}
To prove existence of \(T\), we may assume that \(S\) is reduced
(\SGARefLink{sga3:xv:subsection:2:2}{Proposition 2.2}). By \SGARefLink{sga3:xv:theorem:3-1}{Theorem 3.1}, \SGARefLink{sga3:xv:theorem:3-1:item:ii}{(ii)}\(\Rightarrow\)\SGARefLink{sga3:xv:theorem:3-1:item:i}{(i)}, it is then
enough to show that the set \(U\) of points \(x\in S\) for which there
exists a sub-torus \(T_x\) of \(G_x\) satisfying
\[
  {}_nT_x=M(n)_x
  \qquad\text{for every }n\text{ that is a power of }q
\]
equals \(\operatorname{ens}(S)\). Such a torus is necessarily unique,
and fpqc descent reduces its existence to existence after an extension
of the residue field of \(x\) (Exposé \SGARefLink{sga3:ix:corollary:4:8:item:b}{IX~4.8(b)}).

Since \(S\) is locally noetherian and connected and \(U\) is nonempty
(it contains the point \(s\) of \SGARefLink{sga3:xv:theorem:4-1:item:b}{Theorem 4.1(b)}), it is enough to prove
the following: if \(x,x'\in S\), if \(x\) is a specialization of
\(x'\), and if one of these two points lies in \(U\), then both lie
in \(U\). The usual method (EGA II~7.1.9) reduces this to the case
where \(S\) is the spectrum of a complete discrete valuation ring
with algebraically closed residue field, which was treated in part
(b). This completes the proof of \SGARefLink{sga3:xv:theorem:4-1}{Theorem 4.1}.
